% -*- mode: context; -*-
\chapter{Trigonometrical Identities}
\startitemize[n, 1*broad]
\item $\because A + B + C = \pi \therefore A + B = \pi - C$

  $\Rightarrow \cos(A + B) = \cos(\pi - C) = \cos C \Rightarrow \sin A\sin B - \cos C = \cos A\cos B$

  $\Rightarrow (\sin A\sin B - \cos C)^2 = \cos^2A\cos^2B$

  $\Rightarrow \sin^2A\sin^2B + \cos^2C - 2\sin A\sin B\cos C = (1 - \sin^2A)(1 - \sin^2B)$

  $\Rightarrow \sin^A + \sin^2B + \cos^2C - 1 = 2\sin A\sin B\cos C$

  $\Rightarrow \sin^2A + \sin^2B - \cos^2C = 2\sin A\sin B\cos C$

\item $A + B + C = 180^\circ \Rightarrow \frac{A}{2} + \frac{B}{2} + \frac{C}{2} = 90^\circ \Rightarrow \frac{A}{2} +
  \frac{B}{2} = 90^\circ - \frac{C}{2}$

  $\Rightarrow \cos\left(\frac{A}{2} + \frac{B}{2}\right) = \cos\left(90^\circ - \frac{C}{2}\right)$

  $\Rightarrow \cos\frac{A}{2}\cos\frac{B}{2} - \sin\frac{A}{2}\sin\frac{B}{2} = \sin\frac{C}{2}$

  $\Rightarrow \sin\frac{C}{2} + \sin\frac{A}{2}\sin\frac{B}{2} = \cos\frac{A}{2}\cos\frac{B}{2}$

  $\Rightarrow \left(\sin\frac{C}{2} + \sin\frac{A}{2}\sin\frac{B}{2}\right)^2 = \cos^2\frac{A}{2}\cos^2\frac{B}{2}$

  $\Rightarrow \sin^2\frac{C}{2} + \sin^2\frac{A}{2}\sin^2\frac{B}{2} + 2\sin\frac{A}{2}\sin\frac{B}{2}\sin\frac{C}{2} = \left(1 -
  \cos^2\frac{A}{2}\right)\left(1 - \cos^2\frac{B}{2}\right)$

  $\sin^2\frac{A}{2} + \sin^2\frac{B}{2} + \sin^2\frac{C}{2} = 1 - 2\sin\frac{A}{2}\sin\frac{B}{2}\sin\frac{C}{2}$

\item Let $A + B = C \Rightarrow \cos(A + B) = \cos C$

  $\Rightarrow \sin A\sin B + \cos C = \cos A\cos B \Rightarrow (\sin A\sin B + \cos C)^2 = \cos^2A\cos^2B$

  $\Rightarrow \sin^2A\sin^2B + \cos^2C + 2\sin A\sin B\cos C = (1 - \sin^2A)(1 - \sin^2B)$

  $\Rightarrow \sin^2A + \sin^2B + 2\sin A\sin B\cos C = \sin^2C$

  $\Rightarrow \sin^2A + \sin^2B + 2\sin A\sin B\cos(A + B) = \sin^2(A + B)$

\item Given, $A + B + C = 180^\circ \Rightarrow A + B = 180^\circ - C \Rightarrow \cos(A + B) = -\cos C$

  $\Rightarrow \cos A\cos B + \cos C = \sin A\sin B \Rightarrow (\cos A\cos B + \cos C)^2 = \sin^2A\sin^2B$

  $\Rightarrow \cos^2A\cos^2B + \cos^2C + 2\cos A\cos B\cos C = (1 - \cos^2A)(1 - \cos^2B)$

  $\Rightarrow \cos^2A + \cos^2B + \cos^2C + 2\cos A\cos B\cos C = 1$

\item We have just proved that $\cos^2A + \cos^2B + \cos^2C + 2\cos A\cos B\cos C = 1$

  $\Rightarrow 3 - \sin^2A - \sin^2B - \sin^2C + 2\cos A\cos B\cos C = 1$

  $\Rightarrow \sin^2A + \sin^2B + \sin^2C = 2(1 + \cos A\cos B\cos C)$

\item Given, $A + B + C = 180^\circ \Rightarrow A + B = 180^\circ - C \Rightarrow \cos(A + B) = -\cos C$

  $\Rightarrow \cos A\cos B = \sin A\sin B - \cos C \Rightarrow \cos^2A\cos^2B = \sin^2A\sin^2B + \cos^2C - 2\sin A\sin
  B\cos C$

  $\Rightarrow \cos^2A\cos^2B = (1 - \cos^2A)(1 - \cos^2B) + \cos^2C - 2\sin A\sin B\cos C$

  $\Rightarrow \cos^2A + \cos^2B - \cos^2C = 1 - 2\sin A\sin B\cos C$

\item Given, $A + B + C = 180^\circ \Rightarrow \frac{A + B + C}{2} = 90^\circ$

  $\Rightarrow \cos\left(\frac{A}{2} + \frac{B}{2}\right) = \sin\frac{C}{2}$

  $\Rightarrow \cos\frac{A}{2}\cos\frac{B}{2} - \sin\frac{C}{2} = \sin\frac{A}{2}\sin\frac{B}{2}$

  $\Rightarrow \cos^2\frac{A}{2}\cos^2\frac{B}{2} + \sin^2\frac{C}{2} - 2\cos\frac{A}{2}\cos\frac{B}{2}\sin\frac{C}{2} =
  \left(1 - \cos^2\frac{A}{2}\right)\left(1 - \cos^2\frac{B}{2}\right)$

  $\Rightarrow \cos^2\frac{A}{2} + \cos^2\frac{B}{2} - \cos^2\frac{C}{2} = 2\cos\frac{A}{2}\cos\frac{B}{2}\sin\frac{C}{2}$

\item Given, $A + B + C = 180^\circ \Rightarrow \frac{A + B + C}{2} = 90^\circ$

  $\Rightarrow \cos\left(\frac{A}{2} + \frac{B}{2}\right) = \sin\frac{C}{2}$

  $\Rightarrow \cos\frac{A}{2}\cos\frac{B}{2} = \sin\frac{A}{2}\sin\frac{B}{2} + \sin\frac{C}{2}$

  $\Rightarrow \cos^2\frac{A}{2}\cos^2\frac{B}{2} = \sin^2\frac{A}{2}\sin^2\frac{B}{2} + \sin^2\frac{C}{2} +
  2\sin\frac{A}{2}\sin\frac{B}{2}\sin\frac{C}{2}$

  $\Rightarrow \cos^2\frac{A}{2}\cos^2\frac{B}{2} = \left(1 - \cos^2\frac{A}{2}\right)\left(1 - \cos^2\frac{B}{2}\right) +
  \sin^2\frac{C}{2} + 2\sin\frac{A}{2}\sin\frac{B}{2}\sin\frac{C}{2}$

  $\cos^2\frac{A}{2} + \cos^2\frac{B}{2} + \cos^2\frac{C}{2} = 2 + 2\sin\frac{A}{2}\sin\frac{B}{2}\sin\frac{C}{2}$

\item Given, $A + B + C = \frac{\pi}{2} \Rightarrow A + B = \frac{\pi}{2} - C \Rightarrow \cos(A + B) = \sin C$

  $\Rightarrow \cos A\cos B = \sin A\sin B + \sin C$

  $\Rightarrow \cos^2A\cos^2B = \sin^2A\sin^2B + \sin^2C + 2\sin A\sin B\sin C$

  $\Rightarrow (1 - \sin^2A)(1 - \sin^2B) = \sin^2A\sin^2B + \sin^2C + 2\sin A\sin B\sin C$

  $\Rightarrow \sin^2A + \sin^2B + \sin^2C = 1 - 2\sin A\sin B\sin C$

\item We have just proven that $\sin^2A + \sin^2B + \sin^2C = 1 - 2\sin A\sin B\sin C$ in previous problem.

  $\Rightarrow 1 - \cos^2A + 1 - \cos^2B + 1 - \cos^2C = 1 - 2\sin A\sin B\sin C$

  $\Rightarrow \cos^2A + \cos^2B + \cos^2C = 2 + 2\sin A\sin B\sin C$

\item Givem $A + B + C = 2\pi \Rightarrow A + B = 2\pi - C \Rightarrow \cos(A + B) = \cos C$

  $\Rightarrow \cos A\cos B - \cos C = \sin A\sin B$

  $\Rightarrow \cos^2A\cos^2B + \cos^2C - 2\cos A\cos B\cos C = \sin^2A\sin^2B = (1 - \cos^2A)(1 - \cos^2B)$

  $\Rightarrow \cos^2A + \cos^2B + \cos^2C - 2\cos A\cos B\cos C = 1$

\item Given $A + B = C \Rightarrow \cos(A + B) = \cos C$

  $\Rightarrow \cos A\cos B - \cos C = \sin A\sin B$

  $\Rightarrow \cos^2A\cos^2B + \cos^2C - 2\cos A\cos B\cos C = \sin^2A\sin^2B$

  $\Rightarrow \cos^2A\cos^2B + \cos^2C - 2\cos A\cos B\cos C = (1 - \cos^2A)(1 - \cos^2B)$

  $\Rightarrow \cos^2A + \cos^2B + \cos^2C - 2\cos A\cos B\cos C = 1$

\item Given $A + B = \frac{\pi}{3} \Rightarrow \cos(A + B) = \cos\frac{\pi}{3} = \frac{1}{2}$

  $\Rightarrow \cos A\cos B - \frac{1}{2} = \sin A\sin B$

  $\Rightarrow \cos^2A\cos^2B - \cos A\cos B + \frac{1}{4} = \sin^2A\sin^2B = (1 - \cos^2A)(1 - \cos^2B)$

  $\Rightarrow \cos^2A + \cos^2B - \cos A\cos B = \frac{3}{4}$

\item From problem 12 we have $A + B = C$ and $\cos^2A + \cos^2B + \cos^2C - 2\cos A\cos B\cos C = 1$

  Substituting $C = A + B$ we get $\cos^2A + \cos^2B + \cos^2(A + B) - 2\cos A\cos B\cos(A + B) = 1$

  $\Rightarrow \cos^2B + \cos^2(A + B) - 2\cos A\cos B\cos(A + B) = 1 - \cos^2A = \sin^2A$ which is independent of
  $B$

\item Given $A + B + C = \pi$ and $A + B = 2C \Rightarrow C = \frac{\pi}{3} \Rightarrow A + B = \pi - \frac{pi}{3}$

  $\cos(A + B) = -\cos\frac{\pi}{3}\Rightarrow \cos A\cos B = \sin A\sin B - \frac{1}{2}$

  $\Rightarrow \cos^2A\cos^2B = \sin^2A\sin^2B - \sin A\sin B + \frac{1}{4}$

  $\Rightarrow (1 - \sin^2A)(1 - \sin^2B) = \sin^2A\sin^2B - \sin A\sin B + \frac{1}{4}$

  $\Rightarrow 4(\sin^2A + \sin^2B - \sin A\sin B) = 3$

\item Given $A + B + C = 2\pi \Rightarrow \cos(B + C) = \cos(2\pi - A) = \cos A$

  $\Rightarrow \cos B\cos C - \cos A = \sin B\sin C$

  $\Rightarrow \cos^B\cos^2C + \cos^2A - 2\cos A\cos B\cos C = \sin^2B\sin^2C = (1 - \cos^2B)(1 - \cos^2C)$

  $\Rightarrow \cos^2B + \cos^2C - \sin^2A - 2\cos A\cos B\cos C = 0$

\item Given $A + B + C = 0 \Rightarrow \cos(A + B) = \cos C$

  $\Rightarrow \cos A\cos B - \cos C = \sin A\sin B$

  $\Rightarrow \cos^2A\cos^2B + \cos^2C - 2\cos A\cos B\cos C = \sin^2A\sin^2B = (1 - \cos^2A)(1 - \cos^2B)$

  $\Rightarrow \cos^2A + \cos^2B + \cos^2C = 1 + 2\cos A\cos B\cos C$

\item Putting $A = B - C, B = C - A$ and $C = A - B$ in 17 we can obtain the desired result.

\item Given $A + B + C = \pi,$ we have to prove that $\sin A\cos B\cos C + \sin B\cos C\cos A + \sin C\cos A\cos B= \sin
  A\sin B\sin C$

  Dividing both sides by $\sin A\sin B\sin C,$ we get

  $\cot B\cot C + \cot C\cot A + \cot A\cot B = 1$

  $A + B = \pi - C\Rightarrow \cot(A + B) = -\cot C$

  $\Rightarrow \frac{\cot A\cot B - 1}{\cot A + \cot B} = -\cot C$

  $\Rightarrow \cot B\cot C + \cot C\cot A + \cot A\cot B = 1$

\item Given, $A + B + C = \pi \Rightarrow A + B = \pi - C$

  $\Rightarrow \tan(A + B) = \tan(\pi - C) = -\tan C$

  $\Rightarrow \frac{\tan A + \tan B}{1 - \tan A\tan B} = -\tan C$

  $\Rightarrow \tan A + \tan B + \tan C = \tan A\tan B\tan C$

\item Given $A + B + C = \pi \Rightarrow \frac{A + B}{2} = \frac{\pi - C}{2}$

  $\Rightarrow \tan\frac{A + B}{2} = \tan\frac{\pi - C}{2}$

  $\Rightarrow \frac{\tan\frac{A}{2} + \tan\frac{B}{2}}{1 - \tan\frac{A}{2}\tan\frac{B}{2}} = \cot\frac{C}{2} =
  \frac{1}{\tan\frac{C}{2}}$

  $\Rightarrow \tan\frac{A}{2}\tan\frac{B}{2} + \tan\frac{B}{2}\tan\frac{C}{2} + \tan\frac{C}{2}\tan\frac{A}{2} = 1$

\item Let $B + C - A = \alpha, C + A - B = \beta, A + B - C = \gamma$

  $\alpha + \beta + \gamma = A + B + C = \pi$

  We have just proven that if $A + B + C = \pi$ then $\Rightarrow \tan A + \tan B + \tan C = \tan A\tan B\tan C$

  Thus, substituting we get, $\Rightarrow \tan\alpha + \tan\beta + \tan\gamma = \tan\alpha\tan\beta\tan\gamma$

  $\Rightarrow \tan(B + C - A) + \tan(C + A - B) + \tan(A + B - C) = \tan(B + C - A)\tan(C + A -
  B)\tan(A + B - C)$

\item Given $A + B + C = \pi\Rightarrow  A + B = \pi - C \Rightarrow \cot(A + B) = \cot(\pi - C)$

  $\Rightarrow \frac{\cot A\cot B - 1}{\cot A + \cot B} = -\cot C$

  $\Rightarrow \cot B\cot C + \cot C\cot A + \cot A\cot B = 1$

\item From previosu problem if $A + B + C = \pi$ then $\Rightarrow \cot B\cot C + \cot C\cot A + \cot A\cot B = 1$

  Given $\cot A + \cot B + \cot C = \sqrt{3}$

  $\Rightarrow \cot^2A + \cot^2B + \cot^2C + 2(\cot A\cot B + \cot B\cot C + \cot C\cot A) = 3$

  $\cot^2A + \cot^2B + \cot^2C = 1$

  $2\cot^2A + 2\cot^2B + 2\cot^2C - 2 = 0$

  $2\cot^2A + 2\cot^2B + 2\cot^2C - 2(\cot A\cot B + \cot B\cot C + \cot C\cot A) = 0$

  $(\cot A - \cot B)^2 + (\cot B - \cot C)^2 + (\cot C - \cot A)^2 = 0$

  This is possible only if $\cot A - \cot B = 0$ i.e. $\cot A = \cot B,$ $\cot B - \cot C = 0$ i.e. $\cot
  B = \cot C$ and $\cot C - \cot A = 0$ i.e. $\cot C = \cot A$

  $\therefore \cot A = \cot B = \cot C \Rightarrow A = B = C$

\item $\because A + B + C + D = 2\pi \Rightarrow A + B = 2\pi - C - D$

  $\Rightarrow \tan(A + B) = -\tan(C + D)$

  $\Rightarrow \frac{\tan A + \tan B}{1 - \tan A\tan B} = -\frac{\tan C + \tan D}{1 - \tan C\tan D}$

  $\Rightarrow (\tan A + \tan B)(1 - \tan C\tan D) = -(1 - \tan A\tan B)(\tan C + \tan D)$

  $\Rightarrow \tan A + \tan B + \tan C + \tan D = \tan A\tan B\tan C + \tan A\tan C\tan D + \tan A\tan B\tan D + \tan
  B\tan C\tan D$

  Dividing both sides by $\tan A\tan B\tan C\tan D,$ we get

  $\frac{\tan A + \tan B + \tan C + \tan D}{\tan A\tan B\tan C\tan D} = \frac{1}{\tan A} + \frac{1}{\tan B} + \frac{1}{\tan
    C} + \frac{1}{\tan D}$

  $\Rightarrow \frac{\tan A + \tan B + \tan C + \tan D}{\cot A + \cot B + \cot C + \cot D} = \tan A\tan B\tan C\tan D$

\item Given $A + B + C = \frac{\pi}{2}\Rightarrow A + B = \frac{\pi}{2} - C$

  $\Rightarrow \cot(A + B) = \cot\left(\frac{\pi}{2} - C\right)$

  $\Rightarrow \frac{\cot A\cot B - 1}{\cot A + \cot B} = \tan C = \frac{1}{\cot C}$

  $\Rightarrow \cot A + \cot B + \cot C = \cot A\cot B\cot C$

\item We have just proven in 26 that $\Rightarrow \cot A + \cot B + \cot C = \cot A\cot B\cot C$

  Dividing both sides by $\cot A\cot B\cot C,$ we get

  $\tan A\tan B + \tan B\tan C + \tan C\tan A = 1$

\item Given $A + B + C = \pi \Rightarrow 3(A + B + C) = 3\pi \Rightarrow 3A + 3B = 3\pi - 3C$

  $\Rightarrow \tan(3A + 3B) = \tan(3\pi - 3C) = -\tan3C$

  $\Rightarrow \frac{\tan 3A + \tan 3B}{1 - \tan3A\tan3B} = -\tan3C$

  $\Rightarrow \tan 3A + \tan 3B + \tan 3C = \tan 3A\tan 3B\tan 3C$

\item Given $A + B + C = \pi \Rightarrow \frac{A + B}{2} = \frac{\pi - C}{2}$

  $\Rightarrow \cot\frac{A + B}{2} = \cot\frac{\pi - C}{2}$

  $\Rightarrow \frac{\cot\frac{A}{2}\cot\frac{B}{2} - 1}{\cot\frac{A}{2} + \cot\frac{B}{2}} = \tan\frac{C}{2} =
  \frac{1}{\cot\frac{C}{2}}$

  $\Rightarrow \cot \frac{A}{2} + \cot \frac{B}{2} + \cot \frac{C}{2} = \cot \frac{A}{2}\cot
  \frac{B}{2}\cot \frac{C}{2}$

\item We have to prove that $\frac{\cot A + \cot B}{\tan A + \tan B} + \frac{\cot B + \cot C}{\tan B + \tan C}
  + \frac{\cot C + \cot A}{\tan C + \tan A} = 1$

  Putting $\tan A = \frac{1}{\cot A}, \tan B = \frac{1}{\cot B}, \tan C = \frac{1}{\cot C},$ we get

  $\cot A\cot B + \cot B\cot C + \cot C\cot A = 1$

  We have already proven above in problem 19.

\item Let $A - B = \alpha, B - C = \beta, C - A = \gamma,$ then

  $\alpha + \beta + \gamma = 0$

  $\Rightarrow \tan(\alpha + \beta) = -\tan\gamma$

  $\Rightarrow \frac{\tan\alpha + \tan\beta}{1 - \tan\alpha\tan\beta} = -\tan\gamma$

  $\tan\alpha + \tan\beta + \tan\gamma = \tan\alpha\tan\beta\tan\gamma$

  Substituting back the values, we get

  $\tan(A - B) + \tan(B - C) + \tan(C - A) = \tan(A - B)\tan(B - C)\tan(C - A)$

\item We have already proven in problem 19 that if $A + B + C = 0,$ then

  $\cot A\cot B + \cot B\cot C + \cot C\cot A = 1$

  Let $A = x + y - z, B = z + x - y, C = y + z - x,$ then

  $A + B + C = x + y + z = 0$

  $\Rightarrow \cot A\cot B + \cot B\cot C + \cot C\cot A = 1$

  Substituting back the values, we get

  $\cot(x + y - z)\cot(z + x - y) + \cot(x + y - z)\cot(y + z - x) + \cot(y + z - x)\cot(z + x - y) = 1$

\item Given $A + B + C= n\pi \Rightarrow \tan(A + B) = \tan(n\pi - C) = -\tan C$

  $\Rightarrow \frac{\tan A + \tan B}{1 - \tan A\tan B} = -\tan C$

  $\Rightarrow \tan A + \tan B + \tan C = \tan A\tan B\tan C$

\item L.H.S $= (\sin2A + \sin2B) + \sin2C = 2\sin(A + B)\cos(A - B) + \sin2C$

  $= 2\sin(\pi - C)\cos(A - B) + \sin2C = 2\sin C\cos(A - B) + 2\sin C\cos C$

  $= 2\sin C[\cos(A - B) + \cos\{\pi - (A + B)\}] = 2\sin C[\cos(A - B) - \cos(A + B)]$

  $= 4\sin A\sin B\sin C$

\item L.H.S. $= (\cos A + \cos B) + \cos C - 1 = 2\cos\frac{A + B}{2}\cos\frac{A - B}{2} + \cos C - 1$

  $= 2\cos\left(\frac{\pi}{2} - \frac{C}{2}\right)\cos\frac{A - B}{2} + \cos C - 1$

  $= 2\sin\frac{C}{2}\cos\frac{A - B}{2} + 1 - 2\sin^2\frac{C}{2} - 1$

  $= 2\sin\frac{C}{2}\left[\cos\frac{A - B}{2} - \sin\frac{C}{2}\right]$

  $= 2\sin\frac{C}{2}\left[\cos\frac{A - B}{2} - \sin\left(\frac{\pi}{2} - \frac{A + B}{2}\right)\right]$

  $= 2\sin\frac{C}{2}\left[\cos\frac{A - B}{2} - \cos\frac{A + B}{2}\right]$

  $= 4\sin\frac{A}{2}\sin\frac{B}{2}\sin\frac{C}{2}$

\item We have proven in 34 and 35 that $\sin 2A + \sin 2B + \sin 2C = 4\sin A\sin B\sin C$ and $\cos A + \cos B + \cos
  C - 1 = 4\sin\frac{A}{2}\sin\frac{B}{2}\sin\frac{C}{2}$ respectively. Thus,

  $\frac{\sin 2A + \sin 2B + \sin 2C}{\cos A + \cos B + \cos C - 1} = \frac{4\sin A\sin B\sin
    C}{4\sin\frac{A}{2}\sin\frac{B}{2}\sin\frac{C}{2}}$

  $=
  \frac{4.2\sin\frac{A}{2}\cos\frac{A}{2}.2\sin\frac{B}{2}\cos\frac{B}{2}.2\sin\frac{C}{2}\cos\frac{C}{2}}{4\sin\frac{A}{2}\sin\frac{B}{2}\sin\frac{C}{2}}$

  $= 8\cos\frac{A}{2}\cos\frac{B}{2}\cos\frac{C}{2}$

\item L.H.S. $= \left(\cos \frac{A}{2} + \cos\frac{B}{2}\right) + \cos\frac{C}{2}$

  $= 2\cos\frac{A + B}{4}\cos\frac{A - B}{4} + \sin\frac{\pi - C}{2}$

  $= 2\cos\frac{\pi - C}{4}\cos\frac{A - B}{4} + 2\sin\frac{\pi - C}{4}\cos\frac{\pi - C}{4}$

  $= 2\cos\frac{\pi - C}{4}\left[\cos\frac{A - B}{4} + \cos\left(\frac{\pi}{2} - \frac{\pi - C}{4}\right)\right]$

  $= 2\cos\frac{\pi - C}{4}2\cos\frac{\pi + A + C - B}{8}\cos\frac{\pi + C - A + B}{8}$

  $= 4\cos\frac{\pi - A}{4}\cos\frac{\pi - B}{4}\cos\frac{\pi - C}{4}$

\item L.H.S. $= \left(\sin\frac{A}{2} + \sin \frac{B}{2}\right) + \sin\frac{C}{2}$

  $= 2\sin\frac{A + B}{4}\cos\frac{A - B}{4} + \cos\frac{\pi - C}{2}$

  $= 2\sin\frac{\pi - C}{4}\cos\frac{A - B}{4} + 1 - 2\sin^2\frac{\pi - C}{4}$

  $=1 + 2\sin\frac{\pi - C}{4}\left[\cos\frac{A - B}{4} - \sin\frac{\pi - C}{4}\right]$

  $= 1 + 2\sin\frac{\pi - C}{4}\left[\cos\frac{A - B}{4} - \cos\frac{\pi + C}{4}\right]$

  $= 1 + 2\sin\frac{\pi - C}{4}.2\sin\frac{\pi + A + C - B}{8}\sin\frac{\pi + C - A + B}{8}$

  $= 1 + 4\sin \frac{B + C}{4}\sin \frac{C + A}{4}\sin \frac{A + B}{4}$

\item L.H.S. $= \frac{1 - \cos A}{2} + \frac{1 - \cos B}{2} - \frac{1 - \cos C}{2}$

  $= \frac{1}{2} - \frac{1}{2}[\cos A + \cos B - \cos C]$

  $\cos A + \cos B - \cos C = 2\cos\frac{A + B}{2}\cos\frac{A - B}{2} - \cos C$

  $= 2\sin\frac{C}{2}\cos\frac{A - B}{2} - 1 + 2\sin^2\frac{C}{2}$

  $= -1 + 2\sin\frac{C}{2}\left[\cos\frac{A - B}{2} + \sin\frac{C}{2}\right]$

  $= -1 + 2\sin\frac{C}{2}\left[\cos\frac{A - B}{2} + \cos\frac{A + B}{2}\right]$

  $= -1 + 2\sin\frac{C}{2}.2\cos\frac{A}{2}\cos\frac{B}{2}$

  $= -1 + 4\cos\frac{A}{2}\cos\frac{B}{2}\sin\frac{C}{2}$

  Thus, L.H.S. $= 1 - 2\cos\frac{A}{2}\cos\frac{B}{2}\sin\frac{C}{2}$

\item L.H.S. $= 1 + \cos56^\circ + (\cos58^\circ - \cos66^\circ)$

  $= 2\cos^228^\circ + 2\sin62^\circ\sin4^\circ$

  $=2\cos^228^\circ + 2\cos28^\circ\sin4^\circ$

  $= 2\cos28^\circ[\sin4^\circ + \cos28^\circ]$

  $= 4\cos28^\circ\cos29^\circ\sin33^\circ$

\item Given $A + B + C = \pi,$ we have to prove that $\cos 2A + \cos 2B - \cos 2C = 1 - 4\sin A\sin B\cos C$

  L.H.S. $=\cos 2A + \cos 2B - \cos 2C = \cos 2A + \cos 2B - \cos[2\pi - 2(A + B)]$

  $= 2\cos(A + B)\cos(A - B) - \cos2(A + B) = 2\cos(A + B)\cos(A - B) - 2\cos^2(A + B) + 1$

  $= 1 + 2\cos(A + B)[\cos(A - B) - \cos(A + B)]$

  $= 1 - 4\sin A\sin B\cos C[\because\cos(A + B) = \cos(\pi - C) = -\cos C]$

\item Given $A + B + C = \pi,$ we have to prove that $\sin 2A + \sin 2B - \sin 2C = 4\cos A\cos B\sin C$

  L.H.S. $= \sin 2A + \sin 2B - \sin 2C = 2\sin(A + B)\cos(A - B) - 2\sin C\cos C$

  $[\because \sin(A + B) = \sin(\pi - C) = \sin C, \cos C = \cos[\pi - (A + B)] = -\cos(A + B)]$

  $=2\sin C[\cos(A - B) + \cos(A + B)]$

  $= 4\cos A\cos B\sin C$

\item Given $A + B + C = \pi,$ we have to prove that $\sin A + \sin B + \sin C = 4\cos \frac{A}{2}\cos \frac{B}{2}\cos
  \frac{C}{2}$

  L.H.S. $= \sin A + \sin B + \sin C = 2\sin\frac{A + B}{2}\cos\frac{A - B}{2} + 2\sin\frac{C}{2}\cos\frac{C}{2}$

  $= 2\sin\frac{\pi - C}{2}\cos\frac{A - B}{2} + 2\sin\frac{C}{2}\cos\frac{C}{2}$

  $= 2\cos\frac{C}{2}\cos\frac{A - B}{2} + 2\sin\frac{\pi - A - B}{2}\cos\frac{C}{2}$

  $= 2\cos\frac{C}{2}[\cos\frac{A - B}{2} + \cos\frac{A + B}{2}]$

  $= 4\cos\frac{A}{2}\cos\frac{B}{2}\cos\frac{C}{2}$

\item L.H.S. $= \cos A + \cos B - \cos C = 2\cos\frac{A + B}{2}\cos\frac{A - B}{2} - 1 + 2\sin^2\frac{C}{2}$

  $= 2\cos\left(\frac{\pi - C}{2}\right)\cos\frac{A - B}{2} + 2\sin^2\frac{C}{2} - 1$

  $= 2\sin\frac{C}{2}\left[\cos\frac{A - B}{2} + \cos\left(\frac{[pi}{2} - \frac{C}{2}\right)\right] - 1$

  $= 2\sin\frac{C}{2}\left[\cos\frac{A - B}{2} + \cos\frac{A + B}{2}\right] - 1$

  $= 4\cos \frac{A}{2}\cos \frac{B}{2}\sin \frac{C}{2} - 1$

\item $B + C - A = \pi - A - A = \pi - 2A, C + A - B = \pi - 2B, A + B - C = \pi - 2C$

  $\Rightarrow$ L.H.S. $= \sin 2A + \sin 2B + \sin 2C$

  We have proven in problem 34 that $\sin 2A + \sin 2B + \sin 2C = 4\sin A\sin B\sin C$

  $\therefore \sin(B + C - A) + \sin(C + A - B) + \sin(A + B - C) = 4\sin A\sin B\sin C$

\item L.H.S. $= \frac{\cos A}{\sin B\sin C} + \frac{\cos B}{\sin C\sin A} + \frac{\cos C}{\sin A\sin B} = 2$

  $= \frac{\cos A\sin A + \cos B\sin B + \cos C\sin C}{\sin A\sin B\sin C}$

  $= \frac{\sin 2A + \sin 2B + \sin 2C}{2\sin A\sin B\sin C}$

  We have proven in problem 34 that $\sin 2A + \sin 2B + \sin 2C = 4\sin A\sin B\sin C$

  $\Rightarrow \frac{\sin 2A + \sin 2B + \sin 2C}{2\sin A\sin B\sin C} = 2$

\item Given $A + B + C = \pi,$ we have to prove that $\frac{\sin 2A + \sin 2B + \sin 2C}{\sin A + \sin B + \sin C} =
  8\sin \frac{A}{2}\sin \frac{B}{2}\sin \frac{C}{2}$

  We have proven in problem 34 that $\sin 2A + \sin 2B + \sin 2C = 4\sin A\sin B\sin C$

  We have also proven in problem 43 that $\sin A + \sin B + \sin C = 4\cos \frac{A}{2}\cos \frac{B}{2}\cos\frac{C}{2}$

  Thus, L.H.S. $= \frac{4\sin A\sin B\sin C}{4\cos \frac{A}{2}\cos \frac{B}{2}\cos\frac{C}{2}}$

  $= 8\sin \frac{A}{2}\sin \frac{B}{2}\sin \frac{C}{2}$

\item Given $x + y + z = \frac{\pi}{2},$ we have to prove that $\cos(x - y - z) + \cos(y - z - x) + \cos(z - x - y) -
  4\cos x\cos y\cos z = 0$

  $x - y - z = x - \frac{\pi}{2} + x = 2x - \frac{\pi}{2}$

  Similarly $y - z - x = 2y - \frac{\pi}{2}$ and $z - x - y = 2z - \frac{\pi}{2}$

  $\therefore$ L.H.S. $= \sin 2x + \sin 2y + \sin 2z - 4\cos x\cos y \cos z$

  Now, $\sin 2x + \sin 2y + \sin 2z = 2\sin(x + y)\cos(x - y) + 2\sin z\cos z$

  $= 2\cos z\cos(x - y) + 2\sin\left(\frac{\pi}{2} - x - y\right)\cos z$

  $= 2\cos z[\cos(x - y) + \cos(x + y)]$

  $= 4\cos x\cos y\cos z$

  $\therefore \sin 2x + \sin 2y + \sin 2z - 4\cos x\cos y \cos z = 0$

\item We have to prove that $\sin(x - y) + \sin(y - z) + \sin(z - x) + 4\sin\frac{x - y}{2}\sin\frac{y - z}{2}\sin \frac{z -
  x}{2} = 0$

  Let $x - y = \alpha, y - z = \beta$ and $z - x = \gamma$ then $\alpha + \beta + \gamma = 0$

  The given equation becomes $\sin\alpha + \sin\beta + \sin\gamma +
  4\sin\frac{\alpha}{2}\sin\frac{\beta}{2}\sin\frac{\gamma}{2} = 0$

  Considering $\sin\alpha + \sin\beta + \sin\gamma$

  $= \sin\frac{\alpha + \beta}{2}\cos\frac{\alpha - \beta}{2} + 2\sin\frac{\gamma}{2}\cos\frac{\gamma}{2}$

  $= -\sin\frac{\gamma}{2}\cos\frac{\alpha - \beta}{2} + 2\sin\frac{\gamma}{2}\cos\frac{\alpha + \beta}{2}$

  $= -4\sin\frac{\alpha}{2}\sin\frac{\beta}{2}\sin\frac{\gamma}{2}$

  Thus, $\sin(x - y) + \sin(y - z) + \sin(z - x) + 4\sin\frac{x - y}{2}\sin\frac{y - z}{2}\sin \frac{z - x}{2} = 0$

\item $B + 2C = \pi - A + C, C + 2A = \pi - B + A, A + 2B = \pi - C + B$

  Thus, L.H.S. $= -[\sin(C - A) + \sin(A - B) + \sin(B - C)]$

  Also, note that $A - B + B - C + C - A = 0$ and we have proven in previous problem that $\sin\alpha + \sin\beta +
  \sin\gamma = 4\sin\frac{\alpha}{2}\sin\frac{\beta}{2}\sin\frac{\gamma}{2}$ when $\alpha + \beta + \gamma = 0$

  Thus, $\sin(B + 2C) + \sin(C + 2A) + \sin(A + 2B) = 4\sin\frac{B - C}{2}\sin\frac{C - A}{2}\sin\frac{A - B}{2}$

\item L.H.S. $= \sin\frac{\pi - A}{2} + \sin\frac{\pi - B}{2} + \sin\frac{\pi - C}{2}$

  Following the result of 43 we can say that

  $\sin\frac{\pi - A}{2} + \sin\frac{\pi - B}{2} + \sin\frac{\pi - C}{2} = 4\cos\frac{\pi - A}{4}\cos\frac{\pi -
    B}{4}\cos\frac{\pi - C}{4}$

\item Let $x = \tan A, y = \tan B, z = \tan C$

  Given, $xy + yz + zx = 1$

  $\therefore \tan A\tan B + \tan B\tan C + \tan C\tan A = 1$

  $\Rightarrow \tan C(\tan A + \tan B) = 1 - \tan A\tan B$

  $\Rightarrow \frac{\tan A + \tan B}{1 - \tan A\tan B} = \frac{1}{\tan C} = \cot C$

  $\Rightarrow \tan(A + B) = \tan\left(\frac{\pi}{2} - C\right)$

  $\Rightarrow A + B = \frac{\pi}{2} - C \Rightarrow A + B + C = \frac{\pi}{2}$

  L.H.S. $= \frac{x}{1 - x^2} + \frac{y}{1 - y^2} + \frac{z}{1 - z^2}$

  $= \frac{\tan A}{1 - \tan^2A} + \frac{\tan B}{1 - \tan^2B} + \frac{\tan C}{1 - \tan^2C}$

  $= \frac{1}{2}(\tan 2A + \tan 2B + \tan 2C)$

  We have already proven that if $2A + 2B + 2C = \pi$ then $\tan2A + \tan2B + \tan2C = \tan2A\tan2B\tan2C$

  $\therefore \frac{1}{2}(\tan 2A + \tan 2B + \tan 2C) = \frac{1}{2}\tan2A\tan2B\tan2C$

  $= \frac{1}{2}\frac{2\tan A}{1 - \tan^2A}.\frac{2\tan B}{1 - \tan^2B}.\frac{2\tan C}{1 - \tan^2C}$

  $= \frac{4xyz}{(1 - x^2)(1 - y^2)(1 - z^2)}$

\item Let $x = \tan A, y = \tan B, z = \tan C$

  Now, $x + y + z = xyz$

  $\Rightarrow \tan A + \tan B + \tan C = \tan A\tan B\tan C$

  $\Rightarrow \tan A + \tan B = \tan C(\tan A\tan B - 1)$

  $\Rightarrow \frac{\tan A + \tan B}{1 - \tan A\tan C} = -\tan C = \tan(\pi - C)$

  $\Rightarrow A + B = \pi - C \Rightarrow A + B + C = \pi$

  L.H.S. $= \frac{3x - x^3}{1 - 3x^2} + \frac{3y - y^3}{1 - 3y^2} + \frac{3z - z^3}{1 - 3z^2}$

  $= \frac{3\tan A - \tan^3A}{1 - 3\tan^2A} + \frac{3\tan B - \tan^3B}{1 - 3\tan^2B} + \frac{3\tan C - \tan^3C}{1 -
    3\tan^2C}$

  $= \tan 3A + \tan 3B + \tan 3C$

  Now following like prebious problem

  $\tan3A + \tan3B + \tan3C = \tan3A\tan3B\tan3C$

  $= \frac{3\tan A - \tan^3A}{1 - 3\tan^2A}\frac{3\tan B - \tan^3B}{1 - 3\tan^2B}\frac{3\tan C - \tan^3C}{1 -
    3\tan^2C}$

  $= \frac{3x - x^3}{1 - 3x^2}.\frac{3y - y^3}{1 - 3y^2}.\frac{3z - z^3}{1 - 3z^2}$

\item Given $x + y + z = xyz$, let $x = \tan A, y = \tan B, z = \tan C$

  $\Rightarrow \tan A + \tan B + \tan C = \tan A\tan B\tan C$

  $\Rightarrow \tan A + \tan B = \tan C(\tan A\tan B - 1)$

  $\Rightarrow \frac{\tan A + \tan B}{1 - \tan A\tan C} = -\tan C = \tan(\pi - C)$

  $\Rightarrow A + B = \pi - C \Rightarrow A + B + C = \pi$

  L.H.S. $= \frac{2x}{1 - x^2} + \frac{2y}{1 - y^2} + \frac{2z}{1 - z^2}$

  $= \frac{2\tan A}{1 - \tan^2A} + \frac{2\tan B}{1 - \tan^2B} + \frac{2\tan C}{1 - \tan^2C}$

  $= \tan 2A + \tan 2B + \tan 2C$

  Following like problem 52

  $\tan 2A + \tan 2B + \tan 2C = \tan2A\tan2B\tan2C = \frac{2x}{1 - x^2}.\frac{2y}{1 - y^2}.\frac{2z}{1 - z^2}$

\item Given $x + y + z = xyz$, let $x = \tan A, y = \tan B, z = \tan C$

  $\Rightarrow \tan A + \tan B + \tan C = \tan A\tan B\tan C$

  $\Rightarrow \tan A + \tan B = \tan C(\tan A\tan B - 1)$

  $\Rightarrow \frac{\tan A + \tan B}{1 - \tan A\tan C} = -\tan C = \tan(\pi - C)$

  $\Rightarrow A + B = \pi - C \Rightarrow A + B + C = \pi$

  Given, $x(1 - y^2)(1 - z^2) + y(1 - z^2)(1 - x^2) + z(1 - x^2)(1 - y^2) = 4xyz$

  Dividing both sides with $(1 - x^2)(1 - y^2)(1 - z^2),$ we get

  $\frac{x}{1 - x^2} + \frac{y}{1 - y^2} + \frac{z}{1 - z^2} = \frac{4xyz}{(1 - x^2)(1 - y^2)(1 - z^2)}$

  L.H.S. $= \frac{x}{1 - x^2} + \frac{y}{1 - y^2} + \frac{z}{1 - z^2} = \frac{1}{2}[\tan 2A + \tan 2B + \tan 2C]$

  $= \frac{1}{2}\tan2A\tan2B\tan2C = \frac{4xyz}{(1 - x^2)(1 - y^2)(1 - z^2)}$

\item L.H.S. $= (\cos A + \cos B) + (\cos C + \cos D)$

  $= 2\cos\frac{A + B}{2}\cos\frac{A - B}{2} + 2\cos\frac{C + D}{2}\cos\frac{C - D}{2}$

  $= 2\cos\frac{A + B}{2}\cos\frac{A - B}{2} + 2\cos\left(\pi - \frac{A + B}{2}\right)\cos\frac{C - D}{2}$

  $= 2\cos\frac{A + B}{2}\cos\frac{A - B}{2} - 2\cos\frac{A + B}{2}\cos\frac{C - D}{2}$

  $= 2\cos\frac{A + B}{2}\left[\cos\frac{A - B}{2} - \cos\frac{C - D}{2}\right]$

  $= 2\cos\frac{A + B}{2}.2\sin\frac{A - B + C - D}{4}\sin\frac{C - D - A + B}{4}$

  $= 4\cos\frac{A + B}{2}\sin\frac{A + C - (B + C)}{4}\sin\frac{B + C - (A + D)}{4}$

  $= 4\cos\frac{A + B}{2}\sin\frac{A + C -(2\pi - A - C)}{4}\sin\frac{B + C - (2\pi - B - C)}{4}$

  $= 4\cos\frac{A + B}{2}\sin\frac{A + C - \pi}{2}\sin\frac{B + C - \pi}{2}$

  $= 4\cos\frac{A + B}{2}\cos\frac{B + C}{2}\cos\frac{C + A}{2}$

\item L.H.S. $= \cos^2S + \cos^2(S - A) + \cos^2(S - B) + \cos^2(S - C)$

  $= \frac{1 + \cos 2S}{2} + \frac{1 + \cos(2S - 2A)}{2} + \frac{1 + \cos(2S - 2B)}{2} + \frac{1 + \cos(2S - 2C)}{2}$

  $= \frac{1}{2}[4 + \{\cos2S + \cos(2S - 2A)\} + \{\cos(2S - 2B) + ]\cos(2S - 2C)\}]$

  $= \frac{1}{2}[4 + 2\cos(2S - A)\cos A + 2\cos(2S - B - C)\cos(C - B)]$

  $= \frac{1}{2}[4 + 2\cos(B + C)\cos A + 2\cos A\cos(C - B)]$

  $= \frac{1}{2}[4 + 2\cos A\{\cos(B + C) + \cos(C - B)\}]$

  $= 2 + 2\cos A\cos B\cos C$

\item If $A + B + C = \pi$ then according to problem 21 $\tan\frac{A}{2}\tan\frac{B}{2} +
  \tan\frac{B}{2}\tan\frac{C}{2} + \tan\frac{C}{2}\tan\frac{A}{2} = 1$

  Let $\tan\frac{A}{2} = x, \tan\frac{B}{2} = y, \tan\frac{C}{2} = z$

  $\Rightarrow xy + yz + xz = 1$

  Now, $(x - y)^2 + (y - z)^2 + (z - x)^2 \geq 0$

  $\Rightarrow x^2 + y^2 + z^2 \geq xy + yz + zx$

  $\Rightarrow x^2 + y^2 + z^2 \geq 1$

  $\tan^2\frac{A}{2} + \tan^2\frac{B}{2} + \tan^2\frac{C}{2}\geq 1$

\item We have proven that if $A + B + C = \pi$ then $\tan A + \tan B + \tan C = \tan A\tan B\tan C$

  Thus, L.H.S. $= \tan A\tan B\tan C(\cot A + \cot B + \cot C)$

  $= \tan B\tan C + \tan C\tan A + \tan A\tan B$

  $= \frac{\sin B\sin C}{\cos B\cos C} + \frac{\sin C\sin A}{\cos C\cos A} + \frac{\sin A\sin B}{\cos A\cos B}$

  $= \frac{\cos A\sin B\sin C + \cos B\sin C\sin A + \cos C\sin A\sin B}{\cos A\cos B\cos C}$

  $= \frac{\sin C[\cos A\sin B + \sin A\cos B] + \cos C\sin A\sin B}{\cos A\cos B\cos C}$

  $= \frac{\sin C\sin(A + B) + \cos C\sin A\sin B}{\cos A\cos B\cos C}$

  $= \frac{\sin^2C + \cos C\sin A\sin B}{\cos A\cos B\cos C}$

  $= \frac{1 - \cos^2C + \cos C\sin A\sin B}{\cos A\cos B\cos C}$

  $= \frac{1 + \cos C[\sin A\sin B - \cos C]}{\cos A\cos B\cos C}$

  $= \frac{1 + \cos C[\sin A\sin B - \cos(\pi - A - B)]}{\cos A\cos B\cos C}$

  $= \frac{1 + \cos C[\sin A\sin B + \cos(A + B)]}{\cos A\cos B\cos C}$

  $= \frac{1 + \cos A\cos B\cos C}{\cos A\cos B\cos C} = 1 + \sec A\sec B\sec C$

\item $\cot B\ + \cot C = \frac{\cos B}{\sin B} + \frac{\cos C}{\sin C}$

  $= \frac{\cos B\sin C + \cos C\sin B}{\sin B\sin C} = \frac{\sin(B + C)}{\sin B\sin C}$

  $= \frac{\sin(\pi - A)}{\sin B\sin C} = \frac{\sin A}{\sin B\sin C}$

  Similarly, $\cot C + \cot A = \frac{\sin B}{\sin C\sin A}$ and $\cot A + \cot B = \frac{\sin C}{\sin A + \sin B}$

\item L.H.S. $= \sin^2A\sin B\cos B + \sin^2A \sin C\cos C + \sin^2B\sin A\cos A + \sin^2B\sin C\cos C + \sin^2C\sin A\cos A +
  \sin^2C\sin B\cos B$

  $= (\sin^2A\sin B\cos B + \sin^2B\sin B\cos B) + (\sin^2A\sin C\cos C + \sin^2C\sin A\cos A) + (\sin^2B\sin C\cos C +
  \sin^2C\sin B\cos B)$

  $= \sin A\sin B\sin(A + B) + \sin A\sin C\sin(A + C) + \sin B\sin C\sin(B + C)$

  $= \sin A\sin B\sin C + \sin A\sin C\sin B + \sin B\sin C\sin A$

  $= 3\sin A\sin B\sin C$

\item L.H.S. $= \cos A - \cos B + \cos C - \cos D = 2\sin\frac{A + B}{2}\sin\frac{B - A}{2} + 2\sin\frac{C + D}{2}\sin\frac{D -
  C}{2}$

  $= 2\sin\frac{A + B}{2}\sin\frac{B - A}{2} + 2\sin\frac{2\pi - (A + B)}{2}\sin\frac{D - C}{2}$

  $= 2\sin\frac{A + B}{2}\left[\sin\frac{B - A}{2} + \sin\frac{D - C}{2}\right]$

  $= 2\sin\frac{A + B}{2}.2\sin\frac{B + D - (A + C)}{4}\cos\frac{B + C - (A + D)}{4}$

  $= 4\sin\frac{A + B}{2}\sin\frac{2\pi - 2(A + C)}{4}\cos\frac{2\pi - 2(A + D)}{4}$

  $= 4\sin\frac{A + B}{2}\sin\frac{A + D}{2}\cos \frac{A + C}{2}$

\item Since $A, B, C, D$ are angles of a cyclic quadrilateral $\therefore A + B + C + D = 2\pi, A + C = \pi, B + D = \pi$

  We have proven in problem 56 that $\cos A + \cos B + \cos C + \cos D = 4\cos\frac{A + B}{2}\cos\frac{B +
    C}{2}\cos\frac{C + A}{2}$

  $\cos\frac{A + C}{2} = \cos\frac{\pi}{2} = 0$

  $\therefore \cos A + \cos B + \cos C + \cos D = 0$

\item We know that $(\cot A - \cot B)^2 + (\cot B - \cot C)^2 + (\cot C - \cot A)^2 \geq 0$

  Also, $\because A + B + C = \pi \therefore \cot A\cot B + \cot B\cot C + \cot C\cot A = 1$

  $\therefore 2\cot^2A + 2\cot^2B + 2\cot^2C \geq 2(\cot A\cot B + \cot B\cot C + \cot C\cot A)$

  $\cot^2A + \cot^2B + \cot^2C \geq 1$

\item $\cos\frac{A}{2}\cos\frac{B - C}{2} = \cos\left(\frac{\pi}{2} - \frac{B + C}{2}\right)\cos\frac{B - C}{2}$

  $= \sin\frac{B + C}{2}\cos\frac{B - C}{2} = \frac{1}{2}\left(2\sin\frac{B + C}{2}\cos\frac{B - C}{2}\right)$

  $= \frac{1}{2}(\sin B + \sin C)$

  Similarly $\cos\frac{B}{2}\cos\frac{C - A}{2} = \frac{1}{2}(\sin A + \sin C)$

  and $\cos\frac{C}{2}\cos\frac{A - B}{2} = \frac{1}{2}(\sin A + \sin B)$

  Adding all these we have desired result.

\item $\sin 3A\sin(B - C) = (3\sin A - 4\sin^3A)\sin(B - C)$

  $= 3\sin A\sin(B - C) - 4\sin^2A.\sin A\sin(B - C)$

  $= 3\sin(B + C)\sin(B - C) - 4\sin^2A\sin(B + C)\sin(B - C)[\because B + C = \pi - A \Rightarrow \sin(B + C) = \sin A]$

  $= \frac{3}{2}(\cos2C - \cos2B) - 2\sin^2A(\cos2C - \cos2B)$

  Now, $2\sin^2A(\cos2C - \cos2B) = (1 - \cos 2A)(\cos2C - \cos2B)$

  $= \cos 2C - \cos 2B - \cos2C \cos 2A + \cos 2A\cos 2B$

  Thus, $\sin 3A\sin(B - C) + \sin 3B\sin(C - A) + \sin3C\sin(A - B) = 0$
\stopitemize

\subject{Challenging Problems}
\startitemize[n, 1*broad]
\item Let $P = \cos A\cos B\cos C$. If $P \leq 0$ there is nothing to prove. If $P > 0$ then all the angles are acute
  (at most one angle can be obtuse), so $\cos C > 0$.

  $2P = \cos C[\cos(A - B) + \cos(A + B)] = \cos C[\cos(A - B) - \cos C] \leq \cos C(1 - \cos C)$

  $= \frac{1}{4} - \left(\cos C - \frac{1}{2}\right)^2 \leq \frac{1}{4}$

  $\therefore P \leq \frac{1}{8}$. Equality holds iff $A = B$ and $\cos C = \frac{1}{2},$ i.e. the triangle is
  equilateral.

\item $2\sin\frac{A}{2}\sin\frac{B}{2} = \cos\frac{A - B}{2} - \cos\frac{A + B}{2} \leq 1 - \sin\frac{C}{2}$

  $\therefore 2\sin\frac{A}{2}\sin\frac{B}{2}\sin\frac{C}{2} \leq \sin\frac{C}{2}\left(1 - \sin\frac{C}{2}\right) =
  \frac{1}{4} - \left(\sin\frac{C}{2} - \frac{1}{2}\right)^2 \leq \frac{1}{4}$

  $\therefore 0 < \sin\frac{A}{2}\sin\frac{B}{2}\sin\frac{C}{2} \leq \frac{1}{8}$

  Since $\cos A + \cos B + \cos C = 1 + 4\sin\frac{A}{2}\sin\frac{B}{2}\sin\frac{C}{2},$ we get $1 < \cos A + \cos B +
  \cos C \leq 1 + \frac{4}{8} = \frac{3}{2}$.

\item We have proved that $\sin^2A + \sin^2B + \sin^2C = 2 + 2\cos A\cos B\cos C$.

  By problem 1, $\cos A\cos B\cos C \leq \frac{1}{8}$

  $\therefore \sin^2A + \sin^2B + \sin^2C \leq 2 + \frac{2}{8} = \frac{9}{4}$, equality for an equilateral triangle.

\item For $0 < u, v < 2\pi$ with $u + v < 2\pi$: $\sin u + \sin v = 2\sin\frac{u + v}{2}\cos\frac{u - v}{2} \leq
  2\sin\frac{u + v}{2},$ with equality iff $u = v$.

  $\therefore \sin A + \sin B + \sin C + \sin\frac{\pi}{3} \leq 2\sin\frac{A + B}{2} + 2\sin\frac{C + \frac{\pi}{3}}{2}$

  $\leq 4\sin\frac{A + B + C + \frac{\pi}{3}}{4} = 4\sin\frac{\pi}{3}$

  $\therefore \sin A + \sin B + \sin C \leq 3\sin\frac{\pi}{3} = \frac{3\sqrt{3}}{2}$

  Equality needs $A = B, C = \frac{\pi}{3}$ and $A + B = C + \frac{\pi}{3},$ i.e. $A = B = C = \frac{\pi}{3}$.

\item $\sin A, \sin B, \sin C$ are positive, so by A.M.--G.M. inequality and problem 4,

  $\sin A\sin B\sin C \leq \left(\frac{\sin A + \sin B + \sin C}{3}\right)^3 \leq \left(\frac{\sqrt{3}}{2}\right)^3 =
  \frac{3\sqrt{3}}{8}$

  Equality holds for an equilateral triangle.

\item Let $x = \tan\frac{A}{2}, y = \tan\frac{B}{2}, z = \tan\frac{C}{2}$, all positive. We know $xy + yz + zx = 1$.

  $(x + y + z)^2 - 3(xy + yz + zx) = \frac{1}{2}[(x - y)^2 + (y - z)^2 + (z - x)^2] \geq 0$

  $\therefore (x + y + z)^2 \geq 3 \Rightarrow x + y + z \geq \sqrt{3}$, equality for an equilateral triangle.

\item With $x, y, z$ as in problem 6, $xy, yz, zx$ are positive, so by A.M.--G.M. inequality

  $1 = xy + yz + zx \geq 3(xy.yz.zx)^{1/3} = 3(xyz)^{2/3}$

  $\therefore xyz \leq 3^{-3/2} = \frac{1}{3\sqrt{3}}$

\item We know that $\cot\frac{A}{2} + \cot\frac{B}{2} + \cot\frac{C}{2} = \cot\frac{A}{2}\cot\frac{B}{2}\cot\frac{C}{2}
  = \frac{1}{\tan\frac{A}{2}\tan\frac{B}{2}\tan\frac{C}{2}}$

  By problem 7 this is $\geq 3\sqrt{3}$, so the least value is $3\sqrt{3},$ attained for an equilateral triangle.

\item Let $x = \tan A, y = \tan B, z = \tan C,$ all positive, and $x + y + z = xyz$.

  By A.M.--G.M. inequality, $xyz = x + y + z \geq 3(xyz)^{1/3} \Rightarrow (xyz)^{2/3} \geq 3$

  $\therefore \tan A\tan B\tan C \geq 3\sqrt{3}$, equality for an equilateral triangle.

\item $\tan^2A + \tan^2B + \tan^2C \geq \frac{(\tan A + \tan B + \tan C)^2}{3} = \frac{(\tan A\tan B\tan C)^2}{3}$

  By problem 9, $\tan A\tan B\tan C \geq 3\sqrt{3}$

  $\therefore \tan^2A + \tan^2B + \tan^2C \geq \frac{27}{3} = 9$

\item $(\cot A + \cot B + \cot C)^2 = \cot^2A + \cot^2B + \cot^2C + 2(\cot A\cot B + \cot B\cot C + \cot C\cot A)$

  $= \cot^2A + \cot^2B + \cot^2C + 2 \geq 1 + 2 = 3$ [we have proved $\cot^2A + \cot^2B + \cot^2C \geq 1$]

  Also $\cot A + \cot B + \cot C = \frac{1 + \cos A\cos B\cos C}{\sin A\sin B\sin C} > 0$ (see problem 42) as $\cos
  A\cos B\cos C > -1$.

  $\therefore \cot A + \cot B + \cot C \geq \sqrt{3}$

\item All secants are positive. By A.M.--G.M. inequality and problem 1,

  $\sec A + \sec B + \sec C \geq \frac{3}{(\cos A\cos B\cos C)^{1/3}} \geq \frac{3}{(1/8)^{1/3}} = 6$

\item By A.M.--G.M. inequality and problem 5,

  ${\rm cosec}\,A + {\rm cosec}\,B + {\rm cosec}\,C \geq \frac{3}{(\sin A\sin B\sin C)^{1/3}} \geq
  \frac{3}{\frac{\sqrt{3}}{2}} = 2\sqrt{3}$

\item By A.M.--G.M. inequality and problem 2,

  ${\rm cosec}\,\frac{A}{2} + {\rm cosec}\,\frac{B}{2} + {\rm cosec}\,\frac{C}{2} \geq
  \frac{3}{\left(\sin\frac{A}{2}\sin\frac{B}{2}\sin\frac{C}{2}\right)^{1/3}} \geq \frac{3}{\frac{1}{2}} = 6$

\item For any real $p, q, r$ we have $pq + qr + rp \leq \frac{(p + q + r)^2}{3}$.

  $\therefore \cos A\cos B + \cos B\cos C + \cos C\cos A \leq \frac{(\cos A + \cos B + \cos C)^2}{3}$

  By problem 2, $0 < \cos A + \cos B + \cos C \leq \frac{3}{2}$, so the R.H.S. $\leq \frac{9}{4}.\frac{1}{3} =
  \frac{3}{4}$

\item $\sin A\sin B + \sin B\sin C + \sin C\sin A \leq \frac{(\sin A + \sin B + \sin C)^2}{3} \leq
  \frac{1}{3}.\frac{27}{4} = \frac{9}{4}$ (problem 4)

  Greatest value $= \frac{9}{4}$, attained for an equilateral triangle.

\item $\sin A + \sin B = 2\sin\frac{A + B}{2}\cos\frac{A - B}{2} = 2\cos\frac{C}{2}\cos\frac{A - B}{2} \leq
  2\cos\frac{C}{2}$

  Similarly $\sin B + \sin C \leq 2\cos\frac{A}{2}$ and $\sin C + \sin A \leq 2\cos\frac{B}{2}$

  Adding, $2(\sin A + \sin B + \sin C) \leq 2\left(\cos\frac{A}{2} + \cos\frac{B}{2} + \cos\frac{C}{2}\right)$.
  Equality iff $A = B = C$.

\item $\cos A + \cos B = 2\cos\frac{A + B}{2}\cos\frac{A - B}{2} = 2\sin\frac{C}{2}\cos\frac{A - B}{2} \leq
  2\sin\frac{C}{2}$

  Similarly $\cos B + \cos C \leq 2\sin\frac{A}{2}$ and $\cos C + \cos A \leq 2\sin\frac{B}{2}$

  Adding and dividing by $2$ we get the result. Equality iff $A = B = C$.

\item We have proved that $\sin\frac{A}{2} + \sin\frac{B}{2} + \sin\frac{C}{2} = 1 + 4\sin\frac{B + C}{4}\sin\frac{C +
  A}{4}\sin\frac{A + B}{4}$, and each factor is positive, so the sum is $> 1$.

  For the upper bound, using $\sin u + \sin v \leq 2\sin\frac{u + v}{2}$ $(0 < u, v < \pi)$:

  $\sin\frac{A}{2} + \sin\frac{B}{2} + \sin\frac{C}{2} + \sin\frac{\pi}{6} \leq 2\sin\frac{A + B}{4} + 2\sin\frac{C +
    \frac{\pi}{3}}{4}$

  $\leq 4\sin\frac{A + B + C + \frac{\pi}{3}}{8} = 4\sin\frac{\pi}{6} = 2$

  $\therefore \sin\frac{A}{2} + \sin\frac{B}{2} + \sin\frac{C}{2} \leq 2 - \frac{1}{2} = \frac{3}{2}$

\item We have proved that $\frac{\sin 2A + \sin 2B + \sin 2C}{\sin A + \sin B + \sin C} =
  8\sin\frac{A}{2}\sin\frac{B}{2}\sin\frac{C}{2}$

  By problem 2, R.H.S. $\leq 1$. Since $\sin A + \sin B + \sin C > 0$,

  $\sin 2A + \sin 2B + \sin 2C \leq \sin A + \sin B + \sin C$

\item $2\cos A\cos B = \cos(A - B) + \cos(A + B) \leq 1 - \cos C = 2\sin^2\frac{C}{2}$

  $\therefore \cos A\cos B \leq \sin^2\frac{C}{2}$, and similarly for the other pairs.

  If $\cos A\cos B\cos C \leq 0$ the result is obvious. Otherwise all three cosines are positive and multiplying the
  three inequalities,

  $(\cos A\cos B\cos C)^2 \leq \left(\sin\frac{A}{2}\sin\frac{B}{2}\sin\frac{C}{2}\right)^2$. Taking positive square
  roots gives the result.

\item We have proved that $(\cot A + \cot B)(\cot B + \cot C)(\cot C + \cot A) = {\rm cosec}\,A\,{\rm cosec}\,B\,{\rm
    cosec}\,C$

  $= \frac{1}{\sin A\sin B\sin C} \geq \frac{8}{3\sqrt{3}}$ by problem 5.

\item $\cos 2A + \cos 2B = 2\cos(A + B)\cos(A - B) = -2\cos C\cos(A - B)$

  $\cos 2C = 2\cos^2C - 1 = -2\cos C\cos(A + B) - 1$

  $\therefore \cos 2A + \cos 2B + \cos 2C = -1 - 2\cos C[\cos(A - B) + \cos(A + B)] = -1 - 4\cos A\cos B\cos C$

  $\geq -1 - \frac{4}{8} = -\frac{3}{2}$ by problem 1.

\item Let $s = \sin\frac{A}{2}$. $\cos B + \cos C = 2\sin\frac{A}{2}\cos\frac{B - C}{2} \leq 2s$ and $\cos A = 1 - 2s^2$

  $\therefore 2\cos A + 3(\cos B + \cos C) \leq 2 - 4s^2 + 6s = \frac{17}{4} - \left(2s - \frac{3}{2}\right)^2 \leq
  \frac{17}{4}$

  Greatest value $= \frac{17}{4}$, attained when $B = C$ and $\sin\frac{A}{2} = \frac{3}{4}$, i.e. $\cos A =
  -\frac{1}{8}$, $B = C = \frac{\pi - A}{2}$.

\item $\sin B\sin C = \frac{1}{2}[\cos(B - C) - \cos(B + C)] = \frac{1}{2}[\cos(B - C) + \cos A] \leq \frac{1 + \cos
  A}{2}$

  $\therefore \sin A + \sin B\sin C \leq \frac{1}{2} + \sin A + \frac{1}{2}\cos A \leq \frac{1}{2} + \sqrt{1 +
    \frac{1}{4}} = \frac{1 + \sqrt{5}}{2}$

  Greatest value $= \frac{1 + \sqrt{5}}{2}$, attained when $\tan A = 2$ and $B = C = \frac{\pi - A}{2}$.

\item $\left(-1, \frac{1}{8}\right]$.

  Upper bound is problem 1. Since $\cos^2A + \cos^2B + \cos^2C = 1 - 2\cos A\cos B\cos C$ and each $\cos^2 < 1$, we get
  $1 - 2\cos A\cos B\cos C < 3 \Rightarrow \cos A\cos B\cos C > -1$.

  Taking $B = C = \theta, A = \pi - 2\theta$, the product is $-\cos 2\theta\cos^2\theta$, which varies continuously from
  $-1$ (as $\theta \to 0$) to $\frac{1}{8}$ (at $\theta = \frac{\pi}{3}$), so every value of $\left(-1,
  \frac{1}{8}\right]$ is taken.

\item $\sin A\sin B = \frac{1}{2}[\cos(A - B) - \cos(A + B)] \leq \frac{1}{2}[1 - \cos(A + B)] = \frac{1}{2}(1 - \sin
  C)$

  $\therefore \sin A\sin B\sin C \leq \frac{1}{2}\sin C(1 - \sin C) = \frac{1}{8} - \frac{1}{2}\left(\sin C -
  \frac{1}{2}\right)^2 \leq \frac{1}{8}$

  Equality when $A = B = C = \frac{\pi}{6}$.

\item $\cos 2A + \cos 2B = 2\cos(A + B)\cos(A - B) = 2\sin C\cos(A - B)$

  $\cos 2C = 1 - 2\sin^2C = 1 - 2\sin C\cos(A + B)$

  $\therefore \cos 2A + \cos 2B + \cos 2C = 1 + 2\sin C[\cos(A - B) - \cos(A + B)] = 1 + 4\sin A\sin B\sin C$

  For positive $A, B, C$, by problem 27 this is $\leq 1 + \frac{4}{8} = \frac{3}{2}$.

\item We have proved that $\sin^2A + \sin^2B + \sin^2C = 2 + 2\cos A\cos B\cos C$

  $\therefore 2 = 2 + 2\cos A\cos B\cos C \Rightarrow \cos A\cos B\cos C = 0$

  Hence one of $A, B, C$ is $90^\circ$, i.e. the triangle is right angled.

\item $\sin^2A + \sin^2B = \sin C = \sin(A + B) = \sin A\cos B + \cos A\sin B$

  $\Rightarrow \sin A(\sin A - \cos B) + \sin B(\sin B - \cos A) = 0$

  $\Rightarrow \sin A\left[\sin A - \sin\left(\frac{\pi}{2} - B\right)\right] + \sin B\left[\sin B -
  \sin\left(\frac{\pi}{2} - A\right)\right] = 0$

  If $A + B > \frac{\pi}{2}$ then $A > \frac{\pi}{2} - B$ and $B > \frac{\pi}{2} - A$ (all in $(0, \frac{\pi}{2})$),
  so both brackets are positive and L.H.S. $> 0$. If $A + B < \frac{\pi}{2}$, both brackets are negative.

  $\therefore A + B = \frac{\pi}{2} \Rightarrow C = 90^\circ$

\item Since $\sin A\sin B > 0$ and $\sin C \leq 1,$

  $1 = \cos A\cos B + \sin A\sin B\sin C \leq \cos A\cos B + \sin A\sin B = \cos(A - B) \leq 1$

  So equality holds throughout: $\sin C = 1$ and $\cos(A - B) = 1$

  $\therefore C = 90^\circ, A = B = 45^\circ$

\item $\sin 3A + \sin 3B = 2\sin\frac{3(A + B)}{2}\cos\frac{3(A - B)}{2} = -2\cos\frac{3C}{2}\cos\frac{3(A - B)}{2}$
  $\left[\because \frac{3(A + B)}{2} = \frac{3\pi}{2} - \frac{3C}{2}\right]$

  $\sin 3C = 2\sin\frac{3C}{2}\cos\frac{3C}{2} = -2\cos\frac{3C}{2}\cos\frac{3(A + B)}{2}$

  $\therefore \sin 3A + \sin 3B + \sin 3C = -2\cos\frac{3C}{2}\left[\cos\frac{3(A - B)}{2} + \cos\frac{3(A +
    B)}{2}\right] = -4\cos\frac{3A}{2}\cos\frac{3B}{2}\cos\frac{3C}{2}$

  So one of the cosines is zero, say $\cos\frac{3A}{2} = 0$. As $0 < \frac{3A}{2} < \frac{3\pi}{2},$ $\frac{3A}{2} =
  \frac{\pi}{2} \Rightarrow A = 60^\circ$.

\item $\cos 3A + \cos 3B = 2\cos\frac{3(A + B)}{2}\cos\frac{3(A - B)}{2} = -2\sin\frac{3C}{2}\cos\frac{3(A - B)}{2}$

  $\cos 3C = 1 - 2\sin^2\frac{3C}{2} = 1 + 2\sin\frac{3C}{2}\cos\frac{3(A + B)}{2}$

  $\therefore \cos 3A + \cos 3B + \cos 3C = 1 - 2\sin\frac{3C}{2}\left[\cos\frac{3(A - B)}{2} - \cos\frac{3(A +
    B)}{2}\right] = 1 - 4\sin\frac{3A}{2}\sin\frac{3B}{2}\sin\frac{3C}{2}$

  So one of the sines is zero, say $\sin\frac{3A}{2} = 0$. As $0 < \frac{3A}{2} < \frac{3\pi}{2},$ $\frac{3A}{2} = \pi
  \Rightarrow A = 120^\circ$.

\item $\cot\frac{A}{2} + \cot\frac{C}{2} = \frac{\sin\frac{A + C}{2}}{\sin\frac{A}{2}\sin\frac{C}{2}} =
  \frac{\cos\frac{B}{2}}{\sin\frac{A}{2}\sin\frac{C}{2}} = 2\cot\frac{B}{2}$

  $\Rightarrow \sin\frac{B}{2} = 2\sin\frac{A}{2}\sin\frac{C}{2} \Rightarrow \cos\frac{A + C}{2} =
  2\sin\frac{A}{2}\sin\frac{C}{2} \Rightarrow \cos\frac{A}{2}\cos\frac{C}{2} = 3\sin\frac{A}{2}\sin\frac{C}{2}$

  $\therefore \cos\frac{A - C}{2} = 4\sin\frac{A}{2}\sin\frac{C}{2} = 2\sin\frac{B}{2}$

  $\sin A + \sin C = 2\cos\frac{B}{2}\cos\frac{A - C}{2} = 4\sin\frac{B}{2}\cos\frac{B}{2} = 2\sin B$

\item Let $\tan A = k, \tan B = 2k, \tan C = 3k$. Since $\tan A + \tan B + \tan C = \tan A\tan B\tan C,$

  $6k = 6k^3 \Rightarrow k^2 = 1$. $k = -1$ would make all three angles obtuse, which is impossible.

  $\therefore k = 1$: $A = 45^\circ, B = \tan^{-1}2, C = \tan^{-1}3$.

\item $\cos A = -\cos(B + C) = \sin B\sin C - \cos B\cos C = \cos B\cos C$

  $\Rightarrow \sin B\sin C = 2\cos B\cos C \Rightarrow \tan B\tan C = 2$

  $\tan A = -\tan(B + C) = -\frac{\tan B + \tan C}{1 - \tan B\tan C} = -\frac{\tan B + \tan C}{1 - 2} = \tan B + \tan C$

\item Let $X = 2A, Y = 2B, Z = 2C$ so that $X + Y + Z = 2\pi$.

  $\sin 2X + \sin 2Y = 2\sin(X + Y)\cos(X - Y) = -2\sin Z\cos(X - Y)$

  $\sin 2Z = 2\sin Z\cos Z = 2\sin Z\cos(X + Y)$

  $\therefore$ L.H.S. $= -2\sin Z[\cos(X - Y) - \cos(X + Y)] = -4\sin X\sin Y\sin Z = -4\sin 2A\sin 2B\sin 2C$

\item With $X = 2A, Y = 2B, Z = 2C,$ $X + Y + Z = 2\pi$.

  $\cos 2X + \cos 2Y = 2\cos(X + Y)\cos(X - Y) = 2\cos Z\cos(X - Y)$

  $\cos 2Z = 2\cos^2Z - 1 = 2\cos Z\cos(X + Y) - 1$

  $\therefore$ L.H.S. $= 2\cos Z[\cos(X - Y) + \cos(X + Y)] - 1 = 4\cos 2A\cos 2B\cos 2C - 1$

\item $\sin^3A\cos(B - C) = \sin^2A\sin(B + C)\cos(B - C) = \frac{1}{2}\sin^2A(\sin 2B + \sin 2C)$

  $\therefore$ L.H.S. $= \frac{1}{2}\sum \sin^2A(\sin 2B + \sin 2C) = 3\sin A\sin B\sin C$

  using the identity proved earlier in this chapter.

\item $\cos A\cos(B - C) = -\cos(B + C)\cos(B - C) = -(\cos^2B - \sin^2C) = \sin^2C - \cos^2B$

  Summing cyclically, L.H.S. $= (\sin^2A + \sin^2B + \sin^2C) - (\cos^2A + \cos^2B + \cos^2C)$

  $= 3 - 2(\cos^2A + \cos^2B + \cos^2C) = 3 - 2(1 - 2\cos A\cos B\cos C) = 1 + 4\cos A\cos B\cos C$

\item $\sin^2A\sin(B - C) = \sin A\sin(B + C)\sin(B - C) = \sin A(\sin^2B - \sin^2C)$

  Let $a = \sin A, b = \sin B, c = \sin C$. L.H.S. $= a(b^2 - c^2) + b(c^2 - a^2) + c(a^2 - b^2)$

  This vanishes when $a = b$, $b = c$ or $c = a$, so it equals $k(a - b)(b - c)(c - a)$; comparing the coefficients of
  $bc^2$ gives $k = 1$.

  $\therefore$ L.H.S. $= (\sin A - \sin B)(\sin B - \sin C)(\sin C - \sin A)$

\item $\cot A + \cot B + \cot C = \frac{\cos A\sin B\sin C + \cos B\sin C\sin A + \cos C\sin A\sin B}{\sin A\sin B\sin
  C}$

  $\cos A = -\cos(B + C) = \sin B\sin C - \cos B\cos C \Rightarrow \sin B\sin C = \cos A + \cos B\cos C$

  $\therefore \cos A\sin B\sin C = \cos^2A + \cos A\cos B\cos C$

  Numerator $= \cos^2A + \cos^2B + \cos^2C + 3\cos A\cos B\cos C = 1 + \cos A\cos B\cos C$

\item $\sin B + \sin C - \sin A = 2\cos\frac{A}{2}\cos\frac{B - C}{2} - 2\sin\frac{A}{2}\cos\frac{A}{2} =
  2\cos\frac{A}{2}\left[\cos\frac{B - C}{2} - \cos\frac{B + C}{2}\right]$

  $= 4\cos\frac{A}{2}\sin\frac{B}{2}\sin\frac{C}{2}$, and similarly for the other two factors.

  Also $\sin A + \sin B + \sin C = 4\cos\frac{A}{2}\cos\frac{B}{2}\cos\frac{C}{2}$

  $\therefore$ L.H.S. $= 256\cos^2\frac{A}{2}\cos^2\frac{B}{2}\cos^2\frac{C}{2}\sin^2\frac{A}{2}\sin^2\frac{B}{2}
  \sin^2\frac{C}{2} = 4\left(8\prod\sin\frac{A}{2}\cos\frac{A}{2}\right)^2 = 4\sin^2A\sin^2B\sin^2C$

\item We know $\cot\frac{A}{2} + \cot\frac{B}{2} + \cot\frac{C}{2} = \cot\frac{A}{2}\cot\frac{B}{2}\cot\frac{C}{2} =
  \frac{\cos\frac{A}{2}\cos\frac{B}{2}\cos\frac{C}{2}}{\sin\frac{A}{2}\sin\frac{B}{2}\sin\frac{C}{2}}$

  Now $4\cos\frac{A}{2}\cos\frac{B}{2}\cos\frac{C}{2} = \sin A + \sin B + \sin C$ and
  $4\sin\frac{A}{2}\sin\frac{B}{2}\sin\frac{C}{2} = \cos A + \cos B + \cos C - 1$

  Dividing gives the result.

\item Since $\tan A\tan B\tan C = \tan A + \tan B + \tan C,$ $\frac{\tan A}{\tan B\tan C} = \frac{\tan^2A}{\tan A\tan
  B\tan C} = \frac{\tan^2A}{\tan A + \tan B + \tan C}$

  $\therefore$ L.H.S. $= \frac{\tan^2A + \tan^2B + \tan^2C}{\tan A + \tan B + \tan C} = \frac{(\sum\tan A)^2 - 2\sum\tan
  A\tan B}{\sum\tan A}$

  $= \tan A + \tan B + \tan C - \frac{2(\tan A\tan B + \tan B\tan C + \tan C\tan A)}{\tan A\tan B\tan C}$

  $= \tan A + \tan B + \tan C - 2(\cot C + \cot A + \cot B)$

\item Since $xy + yz + zx = 1,$ $1 + x^2 = x^2 + xy + yz + zx = (x + y)(x + z)$

  $\therefore \frac{x}{1 + x^2} = \frac{x}{(x + y)(x + z)} = \frac{x(y + z)}{(x + y)(y + z)(z + x)}$

  Summing, L.H.S. $= \frac{x(y + z) + y(z + x) + z(x + y)}{(x + y)(y + z)(z + x)} = \frac{2(xy + yz + zx)}{(x + y)(y +
    z)(z + x)} = \frac{2}{(x + y)(y + z)(z + x)}$

\item Put $x = \tan\alpha, y = \tan\beta, z = \tan\gamma$ with $\alpha, \beta, \gamma \in \left(0, \frac{\pi}{2}\right)$.
  Then $1 - xy = z(x + y) > 0$ and $\tan(\alpha + \beta) = \frac{x + y}{1 - xy} = \frac{1}{z} = \cot\gamma$

  $\therefore \alpha + \beta + \gamma = \frac{\pi}{2}$, i.e. $2\alpha + 2\beta + 2\gamma = \pi$.

  $\frac{1}{1 + x^2} = \cos^2\alpha$ and $\cos^2\alpha + \cos^2\beta + \cos^2\gamma = 2 + 2\sin\alpha\sin\beta\sin\gamma$
  (proved earlier for half angles of a triangle)

  $\leq 2 + \frac{2}{8} = \frac{9}{4}$ by problem 2.

\item Put $x = \tan A, y = \tan B, z = \tan C$ with $A, B, C \in \left(0, \frac{\pi}{2}\right)$. From $x + y + z =
  xyz,$ $z(xy - 1) = x + y > 0 \Rightarrow xy > 1$.

  $\tan(A + B) = \frac{x + y}{1 - xy} = -z = -\tan C$ and $\frac{\pi}{2} < A + B < \pi$ $\therefore A + B + C = \pi$

  $\frac{1}{\sqrt{1 + x^2}} = \cos A$ etc., so L.H.S. $= \cos A + \cos B + \cos C \leq \frac{3}{2}$ by problem 2.

\item Let $A = \frac{\pi}{7}, B = \frac{2\pi}{7}, C = \frac{4\pi}{7}$; then $A + B + C = \pi$ and $\cos^2A + \cos^2B +
  \cos^2C = 1 - 2\cos A\cos B\cos C$

  $\cos A\cos B\cos C = \frac{8\sin A\cos A\cos 2A\cos 4A}{8\sin A} = \frac{\sin 8A}{8\sin A} = \frac{\sin(\pi +
    \frac{\pi}{7})}{8\sin\frac{\pi}{7}} = -\frac{1}{8}$

  $\therefore \cos^2\frac{\pi}{7} + \cos^2\frac{2\pi}{7} + \cos^2\frac{4\pi}{7} = 1 + \frac{2}{8} = \frac{5}{4}$

\item $\sin^3A = \frac{3\sin A - \sin 3A}{4}$

  $\therefore \sin^3A + \sin^3B + \sin^3C = \frac{3}{4}(\sin A + \sin B + \sin C) - \frac{1}{4}(\sin 3A + \sin 3B + \sin
  3C)$

  $= \frac{3}{4}.4\cos\frac{A}{2}\cos\frac{B}{2}\cos\frac{C}{2} - \frac{1}{4}\left(-4\cos\frac{3A}{2}\cos\frac{3B}{2}
  \cos\frac{3C}{2}\right)$ (see problem 32)

  $= 3\cos\frac{A}{2}\cos\frac{B}{2}\cos\frac{C}{2} + \cos\frac{3A}{2}\cos\frac{3B}{2}\cos\frac{3C}{2}$
\stopitemize
